%% toc: P15 157 157.3,157.4,157.5,157.6
\chapter{Lebesgue Integral Notation and the Change of Variables Formula}
\label{sec:p15-w02c-p15-001-integral-notation}

Two pieces of notation run through the rest of this part. The first writes
an integral against Lebesgue measure with a variable of integration, as
\(\int f(x)\,dx\) or \(\int_a^b f(x)\,dx\); later chapters use it for
marginal densities, conditional densities and the density calculations of
\cref{sec:p15-convolution-gamma-normal} and \cref{sec:p15-radon-nikodym}.
The second is the name \emph{change of variables formula} for the identity
\(\mathbb E[f(X)]=\int_S f\,d\mathbb P_X\) of
\cref{thm:p15-expectation-law-integral}, together with the reading of
that identity which explains the name. This chapter fixes both, proves the
facts about endpoints, additivity and limits that the interval notation
silently uses, and shows the formula at work on the uniform variable.

\emph{Inputs.} The integral of
\cref{def:p15-nonnegative-integral,def:p15-integrable-functions}; Borel
Lebesgue measure \(\lambda\) on \(\mathbb R\)
(\cref{thm:p15-lebesgue-unit-interval}) and the product measure
\(\lambda^d\) on \(\mathcal B(\mathbb R^d)\)
(\cref{thm:p15-finite-product-measures}, with \(\lambda^1=\lambda\) as
in \cref{def:p15-probability-density}); the law \(\mathbb P_X\) of a
random element (\cref{def:p15-pushforward-law}); the uniform
probability space \((\Omega_U,\mathcal F_U,\mathbb U)\) with its identity
variable \(U\); and, for the identification of \(\lambda^d\), the
uniqueness of a measure that agrees on a pi-system with a finite exhaustion
(\cref{thm:p15-w02c-p15-002-sigma-finite-uniqueness}).
\emph{Conclusions.} The meaning of the interval and \(dx\) notations, of
\(\int_S f(y)\,\mu(dy)\), of \(\int f(x)\,dG(x)\) for a Stieltjes
measure and of \(\sum_if(i)\) as an integral against counting measure,
with proofs that these notations name well-defined integrals; the fact that rectangle volumes determine
\(\lambda^d\), so that Durrett's Lebesgue measure on \(\mathbb R^d\) and
\(\lambda^d\) coincide; independence of the interval integral from the
endpoints included, and its additivity in the interval; the relation
between the interval integral and its limits as an endpoint tends to
infinity, with a function for which that limit exists although the
Lebesgue integral over the half-line is undefined; the change of variables
formula in Durrett's two forms, with his motive for it; Durrett's
four-step method of proof, stated as a procedure and proved as a
theorem; and the identity
\(\mathbb E[h(U)]=\int_0^1h(x)\,dx\), which computes
\(\int_0^1x\,dx=1/2\) and \(\int_0^1x^2\,dx=1/3\) without any
appeal to a fundamental theorem of calculus. Three further results use
this notation: Durrett's direct proof of the Cauchy--Schwarz inequality
from a nonnegative quadratic; his squared form of Chebyshev's inequality,
with the competing uses of the names Chebyshev and Markov; and the values
of the interval integrals of affine functions and of the reciprocal
function that the density examples of \cref{sec:p15-radon-nikodym}
quote. The chapter closes with the comparison of absolute moments of
different orders and with the mean and variance of the Bernoulli law
(\cref{sec:p15-w02c-p15-001-lower-moments}), and with the role of each
hypothesis of Jensen's inequality: the total mass one, the integrability
of \(\varphi(X)\) and the integrability of \(X\)
(\cref{sec:p15-w02c-p15-001-jensen-hypotheses}).

\emph{Source.} Durrett's author version of \emph{Probability: Theory and
Examples}, dated 11 January 2019. After proving the basic properties of the
integral he lists four notations for special cases
\cite[sec.~1.4, pp.~23--24]{DurrettPTEAuthor2019}. He states
Theorem~1.6.9 under the title ``Change of variables formula'' and explains
the name in the remark that follows it
\cite[thm.~1.6.9, p.~32]{DurrettPTEAuthor2019}; his proof runs through
indicators, simple functions, nonnegative functions and integrable
functions \cite[pp.~32--33]{DurrettPTEAuthor2019}. Durrett writes
\(\mathcal R^d\) for the Borel subsets of \(\mathbb R^d\), dropping the
superscript when \(d=1\) \cite[p.~3]{DurrettPTEAuthor2019}; this handbook
writes \(\mathcal B(\mathbb R^d)\) (\cref{def:p15-borel-sets}).

\section{Integrals against Lebesgue measure}
\label{sec:p15-w02c-p15-001-dx-notation}

\begin{hbconvention}{def:p15-w02c-p15-001-dx-notation}{Writing integrals against Lebesgue measure}
Let \(d\geq1\) and let \(f:\mathbb R^d\to[-\infty,\infty]\) be Borel
measurable. For a Borel set \(E\), the expression
\(\int_E f\,d\lambda^d\) means \(\int f\mathbf1_E\,d\lambda^d\), as in
\cref{def:p15-nonnegative-integral,def:p15-integrable-functions}. It has a
value exactly when that integral exists: always when \(f\geq0\) on
\(E\), and otherwise when the integrals of \(f^+\mathbf1_E\) and
\(f^-\mathbf1_E\) are not both infinite. The symbol \(\int\) without a
subscript integrates over the whole space.
\begin{enumerate}
\item \emph{Lebesgue measure on
      \(\mathbb R^d\).} In Durrett's case~(a), the measure space
      \((\Omega,\mathcal F,\mu)\) is
      \((\mathbb R^d,\mathcal R^d,\lambda)\), and \(\int f(x)\,dx\) denotes \(\int f\,d\lambda\)
      \cite[sec.~1.4, p.~23]{DurrettPTEAuthor2019}. In this handbook
      the measure space is \((\mathbb R^d,\mathcal B(\mathbb R^d),\lambda^d)\),
      so
      \[
       \int f(x)\,dx\defeq\int f\,d\lambda^d,\qquad
       \int_E f(x)\,dx\defeq\int_E f\,d\lambda^d .
      \]
\item \emph{A closed bounded
      interval.} Durrett's case~(b): for
      \((\Omega,\mathcal F,\mu)=(\mathbb R,\mathcal R,\lambda)\) and
      \(E=[a,b]\), \(\int_a^b f(x)\,dx\) denotes \(\int_E f\,d\lambda\)
      \cite[sec.~1.4, p.~23]{DurrettPTEAuthor2019}. Thus, for real
      numbers \(a\leq b\),
      \[
       \int_a^b f(x)\,dx\defeq\int_{[a,b]}f\,d\lambda .
      \]
      The two limits of integration name the interval; they are not an
      orientation, and no value is assigned here to \(\int_a^b\) with
      \(a>b\). Case~(b) is a Lebesgue integral over a bounded interval.
      An expression written with an infinite limit is not covered by it;
      where a later chapter writes one, that chapter states what it
      means (see \cref{thm:p15-w02c-p15-001-interval-limits} for the
      difference this makes).
\item \emph{The variable of
      integration.} The letter \(x\) is bound by \(dx\): \(\int f(t)\,dt\)
      and \(\int f(x)\,dx\) are the same number. When the integrand has
      further arguments, the bound variable selects the section that is
      integrated. For a Borel \(f\) on \(\mathbb R\times\mathbb R\) and
      fixed \(y\), \(\int f(x,y)\,dx\) is the integral against
      \(\lambda\) of the section \(f^y(x)=f(x,y)\) of
      \cref{def:p15-product-events}; the section is Borel by
      \cref{thm:p15-product-sections}. This is how the formula for
      \(m(y)\) in \cref{thm:p15-conditional-density-formulas} reads.
\item \emph{A named measure with a
      named variable.} For a measure \(\mu\) on \((S,\mathcal S)\),
      Durrett writes \(\int_S f(y)\,\mu(dy)\) for \(\int_S f\,d\mu\)
      \cite[thm.~1.6.9, p.~32]{DurrettPTEAuthor2019}, and
      \(\int_S f(y)\,d\mu(y)\) is used with the same meaning.
\item \emph{A Stieltjes measure.} Durrett's case~(c): when
      \((\Omega,\mathcal F,\mu)=(\mathbb R,\mathcal R,\mu)\) with
      \(\mu((a,b])=G(b)-G(a)\) for \(a<b\), he writes
      \(\int f(x)\,dG(x)\) for \(\int f\,d\mu\)
      \cite[sec.~1.4, p.~23]{DurrettPTEAuthor2019}. Here \(G\) is a
      real function on \(\mathbb R\). By
      \cref{thm:p15-w02c-p15-001-stieltjes-notation}, such a \(\mu\)
      exists exactly when \(G\) is a Stieltjes measure function in the
      sense of \cref{def:p15-stieltjes-intervals}, and it is then the
      measure \(\mu_G\) of \cref{thm:p15-stieltjes-construction}; so
      \[
       \int f(x)\,dG(x)\defeq\int f\,d\mu_G ,\qquad
       \int_E f(x)\,dG(x)\defeq\int_E f\,d\mu_G ,
      \]
      with the existence conditions stated at the start of this
      convention. The symbol \(dG(x)\) does not assume that \(G\) has a
      derivative, and the letter \(x\) is bound as in item~(3). When
      \(G\) is the cumulative distribution function of a probability law
      \(\nu\), \(\mu_G=\nu\); this is the reading of
      \cref{def:p15-convolution-law}.
\item \emph{Counting measure.} Durrett's case~(d): when \(\Omega\) is a
      countable set, \(\mathcal F\) consists of all its subsets and
      \(\mu\) is counting measure, he writes \(\sum_i f(i)\) for
      \(\int f\,d\mu\) \cite[sec.~1.4, p.~24]{DurrettPTEAuthor2019}.
      \emph{Counting measure} on a set \(I\) with all its subsets is
      \(\#(A)\defeq\) the number of elements of \(A\), and \(\infty\) when
      \(A\) is infinite; it is countably additive, because a disjoint
      union of sets is finite exactly when finitely many of them are
      nonempty and all are finite, and then its size is the sum of their
      sizes. For a countable \(I\),
      \[
       \sum_{i\in I}f(i)\defeq\int f\,d\#
      \]
      whenever the integral exists.
      \Cref{thm:p15-w02c-p15-001-counting-sums} shows that this agrees
      with the sum of a nonnegative series in every order, and with the
      sum of an absolutely convergent series for integrable \(f\). Durrett
      states that he records this case mainly to indicate that results
      for sums follow from results for integrals
      \cite[sec.~1.4, p.~24]{DurrettPTEAuthor2019}; part~(4) of that
      lemma is an instance.
\end{enumerate}
\index{integral!notation \(dx\)}\index{Lebesgue measure!integral notation}
\index{integral!notation \(dG(x)\)}\index{counting measure!sums as integrals}
\index{Stieltjes measure!integral notation}
\end{hbconvention}

Durrett obtains Lebesgue measure on \(\mathbb R^d\) from his Example~1.1.12
with every one-dimensional factor equal to the identity
\cite[ex.~1.1.12, p.~8]{DurrettPTEAuthor2019}. The measure \(\lambda^d\)
used here is the product measure of \cref{thm:p15-finite-product-measures};
both give each finite rectangle \(\prod_{i=1}^d(a_i,b_i]\) its volume
\(\prod_{i=1}^d(b_i-a_i)\). The next lemma shows that this property
determines the measure on all Borel sets, so the two constructions give the
same measure and case~(a) has the same meaning in both.

\begin{hblemma}{thm:p15-w02c-p15-001-lebesgue-rectangles}{proved=annex:p15-w02c-p15-001-lebesgue-rectangles}{Lebesgue measure on \(\mathbb R^d\) is fixed by rectangle volumes}
\par\nobreak
Let \(d\geq1\) and let \(\nu\) be a measure on
\((\mathbb R^d,\mathcal B(\mathbb R^d))\) such that
\[
 \nu\left(\prod_{i=1}^d(a_i,b_i]\right)=\prod_{i=1}^d(b_i-a_i)
 \qquad(a_i<b_i\text{ real}).
\]
Then \(\nu=\lambda^d\).
\end{hblemma}

Durrett's case~(c) names a measure by its values on half-open intervals
alone. The next lemma shows that those values fix the measure and force
\(G\) to be a Stieltjes measure function, and it computes three instances.

\begin{hblemma}{thm:p15-w02c-p15-001-stieltjes-notation}{proved=annex:p15-w02c-p15-001-stieltjes-notation}{The measure named by \(dG(x)\)}
\par\nobreak
Let \(G:\mathbb R\to\mathbb R\).
\begin{enumerate}
\item If \(\mu\) is a measure on \((\mathbb R,\mathcal B(\mathbb R))\) with
      \(\mu((a,b])=G(b)-G(a)\) for all real \(a<b\), then \(G\) is
      nondecreasing and right-continuous, and \(\mu=\mu_G\), the measure
      of \cref{thm:p15-stieltjes-construction}.
\item Conversely, if \(G\) is a Stieltjes measure function, then
      \(\mu_G((a,b])=G(b)-G(a)\) for all real \(a<b\).
\item If \(G=F_X\) is the cumulative distribution function of a real
      random variable \(X\) (\cref{def:p15-cumulative-distribution}),
      then \(\mu_G=\mathbb P_X\), so
      \(\int f(x)\,dF_X(x)=\int f\,d\mathbb P_X\).
\item If \(G(x)=x\), then \(\mu_G=\lambda\), so
      \(\int f(x)\,dG(x)=\int f(x)\,dx\).
\item If \(G=\mathbf1_{[0,\infty)}\), then \(\mu_G=\delta_0\), and
      \(\int f(x)\,dG(x)=f(0)\) for every Borel
      \(f:\mathbb R\to[0,\infty]\) and every real Borel \(f\); in the
      second case \(f\) is \(\mu_G\)-integrable.
\end{enumerate}
\end{hblemma}

\begin{hblemma}{thm:p15-w02c-p15-001-counting-sums}{proved=annex:p15-w02c-p15-001-counting-sums}{Integrals against counting measure are sums}
\par\nobreak
\begin{enumerate}
\item Let \((\Omega,\mathcal F,\mu)\) be a measure space, let
      \(s\in\Omega\) with \(\{s\}\in\mathcal F\), and let
      \(f:\Omega\to[0,\infty]\) be measurable. Then
      \(\int_{\{s\}}f\,d\mu=f(s)\,\mu(\{s\})\), with the products of
      \cref{def:p15-simple-integral}.
\end{enumerate}
Now let \(I\) be a nonempty countable set and \(\#\) counting measure on
all subsets of \(I\). Call a sequence \((i_k)\) an \emph{enumeration} of
\(I\) when it lists every element of \(I\) exactly once (a finite list
when \(I\) is finite).
\begin{enumerate}\setcounter{enumi}{1}
\item For \(f:I\to[0,\infty]\),
      \[
       \int f\,d\#=\sup\Big\{\sum_{i\in F}f(i):F\subseteq I\text{ finite}\Big\}
       =\lim_{N\to\infty}\sum_{k\leq N}f(i_k)
      \]
      for every enumeration \((i_k)\) of \(I\), in \([0,\infty]\). In
      particular the value does not depend on the order of summation.
\item A function \(f:I\to\mathbb R\) is \(\#\)-integrable exactly when
      \(\sum_{i\in I}|f(i)|<\infty\). Then, for every enumeration
      \((i_k)\), the series \(\sum_kf(i_k)\) converges absolutely and
      \(\int f\,d\#=\sum_kf(i_k)\).
\item If \(f_n:I\to[0,\infty]\) satisfy \(f_n\leq f_{n+1}\) and
      \(f_n\to f\) pointwise, then
      \(\sum_{i\in I}f_n(i)\uparrow\sum_{i\in I}f(i)\).
\end{enumerate}
\end{hblemma}

Part~(4) is the monotone convergence theorem for series, read off from
\cref{thm:p15-monotone-integral} with \(\mu=\#\). This is the use of
case~(d) that Durrett announces: a theorem about integrals, applied to
counting measure, becomes a theorem about sums.

\begin{hbexample}{def:p15-w02c-p15-001-stieltjes-counting-examples}{A point mass, a geometric sum and a series without an integral}
\begin{enumerate}
\item With \(G=\mathbf1_{[0,\infty)}\), part~(5) of
      \cref{thm:p15-w02c-p15-001-stieltjes-notation} gives
      \(\int(x^2+1)\,dG(x)=1\). The same integral against \(G(x)=x\) is
      \(\int(x^2+1)\,dx=\infty\), since \(x^2+1\geq1\) and
      \(\lambda(\mathbb R)=\infty\). The notation \(dG(x)\) carries the
      measure, and the integrand alone does not determine the value.
\item On \(I=\{1,2,\ldots\}\), let \(f(i)=2^{-i}\). For the enumeration
      \(i_k=k\), the partial sums are \(\sum_{k\leq N}2^{-k}=1-2^{-N}\)
      by the finite geometric identity of
      \cref{sec:p15-measurable-probability}, so part~(2) of
      \cref{thm:p15-w02c-p15-001-counting-sums} gives \(\int f\,d\#=1\).
\item On the same \(I\), let \(f(i)=(-1)^{i+1}/i\). Then
      \(\int f^-\,d\#=\sum_{k\geq1}1/(2k)\geq\frac12\sum_{k\geq1}1/(k+1)\),
      which is infinite because the last series diverges
      (\cref{def:p15-convergence-failures}); and
      \(\int f^+\,d\#=\sum_{k\geq1}1/(2k-1)\geq\sum_{k\geq1}1/(2k)=\infty\).
      So \(\int f\,d\#\) is undefined. Yet the partial sums
      \(\sum_{k\leq N}f(k)\), taken in the order \(1,2,3,\ldots\), are the
      numbers \(S_N\) of part~(3) of
      \cref{thm:p15-w02c-p15-001-interval-limits}, and they converge to
      the limit \(L\) there. The sum notation of case~(d) names the
      integral only, not such an ordered limit; this series is the
      discrete counterpart of the function \(g\) of that theorem.
\end{enumerate}
\end{hbexample}

\begin{hblemma}{thm:p15-w02c-p15-001-interval-integrals}{proved=annex:p15-w02c-p15-001-interval-integrals}{Endpoints and additivity of interval integrals}
\par\nobreak
Let \(f:\mathbb R\to[-\infty,\infty]\) be Borel and let \(a\leq b\) be
real.
\begin{enumerate}
\item If \(f\geq0\) on \([a,b]\), or \(\int_a^b|f(x)|\,dx<\infty\), then
      \[
       \int_a^b f(x)\,dx=\int_{(a,b)}f\,d\lambda=\int_{(a,b]}f\,d\lambda
       =\int_{[a,b)}f\,d\lambda,
      \]
      and in the second case all four integrals are finite.
\item Under the same hypothesis, for every \(c\) with \(a\leq c\leq b\),
      \[
       \int_a^b f(x)\,dx=\int_a^c f(x)\,dx+\int_c^b f(x)\,dx ,
      \]
      a sum of nonnegative extended numbers in the first case and of
      finite numbers in the second.
\item For a constant \(k\in\mathbb R\), \(\int_a^b k\,dx=k(b-a)\).
\end{enumerate}
\end{hblemma}

\begin{hbtheorem}{thm:p15-w02c-p15-001-interval-limits}{proved=annex:p15-w02c-p15-001-interval-limits}{Interval integrals, their limits and the half-line integral}
\par\nobreak
\begin{enumerate}
\item If \(f:\mathbb R\to[0,\infty]\) is Borel, then
      \(b\mapsto\int_0^b f(x)\,dx\) is nondecreasing on \([0,\infty)\)
      and
      \[
       \lim_{b\to\infty}\int_0^b f(x)\,dx=\int_{[0,\infty)}f\,d\lambda,
      \]
      both sides possibly infinite.
\item If \(f\) is real Borel and \(\int_{[0,\infty)}|f|\,d\lambda<\infty\),
      then each \(\int_0^b f(x)\,dx\) is finite and converges to the finite
      number \(\int_{[0,\infty)}f\,d\lambda\) as \(b\to\infty\).
\item Let
      \[
       c_n\defeq\frac{(-1)^n}{n+1},\qquad
       g\defeq\sum_{n=0}^\infty c_n\mathbf1_{[n,n+1)},
      \]
      so that \(g=c_n\) on \([n,n+1)\) and \(g=0\) on \((-\infty,0)\).
      Then \(g\) is Borel and bounded, and for real \(b\geq0\) with
      \(N=\lfloor b\rfloor\),
      \[
       \int_0^b g(x)\,dx=S_N+(b-N)c_N,\qquad
       S_N\defeq\sum_{n=0}^{N-1}c_n .
      \]
      The limit \(L\defeq\lim_{b\to\infty}\int_0^b g(x)\,dx\) exists and
      \(1/2\leq L\leq5/6\). Nevertheless
      \(\int_{[0,\infty)}g^+\,d\lambda=\int_{[0,\infty)}g^-\,d\lambda=\infty\),
      so \(\int_{[0,\infty)}g\,d\lambda\) is undefined.
\end{enumerate}
\end{hbtheorem}

\begin{hbwarning}{def:p15-w02c-p15-001-limit-not-integral}{A limit of interval integrals is a different object}
For nonnegative and for absolutely integrable functions, letting an
endpoint of a bounded interval tend to infinity recovers the Lebesgue
integral over the half-line, by parts~(1) and~(2) of
\cref{thm:p15-w02c-p15-001-interval-limits}. For a signed function whose
positive and negative parts both have infinite integral, the limit of the
interval integrals can exist while the half-line integral is undefined, as
the function \(g\) of part~(3) shows. The limit then depends on the order
in which the line is exhausted, which is exactly the cancellation that
\cref{def:p15-integrable-functions} refuses to perform. Case~(b) of
\cref{def:p15-w02c-p15-001-dx-notation} therefore never names such a limit.
\end{hbwarning}

\section{The change of variables formula}
\label{sec:p15-w02c-p15-001-change-of-variables}

\begin{hbconvention}{def:p15-w02c-p15-001-change-of-variables-name}{The name of the law-integral identity}
For a random element \(X\) of \((S,\mathcal S)\) on
\((\Omega,\mathcal F,\mathbb P)\), the identity
\(\mathbb E[f(X)]=\int_S f\,d\mathbb P_X\) of
\cref{thm:p15-expectation-law-integral} is Durrett's Theorem~1.6.9, which he
titles the \emph{change of variables formula}
\cite[thm.~1.6.9, p.~32]{DurrettPTEAuthor2019}. He introduces it with a
motive: an integral over the abstract space \((\Omega,\mathcal F,\mathbb P)\)
serves the theory, but a computation has to be moved to a space on which
calculus can be done, and in most of his applications that space is
\(S=\mathbb R^d\) \cite[sec.~1.6.3, p.~32]{DurrettPTEAuthor2019}. The
formula is the vehicle for that move. In the remark after the
theorem he explains the name by renaming the random element and its law:
write \(h\) for \(X\) and \(\mathbb P\circ h^{-1}\) for its distribution
\(\mu\), where
\[
 (\mathbb P\circ h^{-1})(B)\defeq\mathbb P(h^{-1}(B))\qquad(B\in\mathcal S),
\]
and the formula becomes
\[
 \int_\Omega f(h(\omega))\,d\mathbb P(\omega)
 =\int_S f(y)\,d(\mathbb P\circ h^{-1})(y)
\]
\cite[thm.~1.6.9 and remark, p.~32]{DurrettPTEAuthor2019}.
Read this way, an integral over the sample space with integration variable
\(\omega\) becomes an integral over the value space with integration
variable \(y=h(\omega)\): the substitution \(y=h(\omega)\) is made in
the integrand, and the measure is carried along by \(h\) to the image
measure \(\mathbb P\circ h^{-1}\). The composite \(\mathbb P\circ h^{-1}\)
applies the inverse-image map \(B\mapsto h^{-1}(B)\) of
\cref{def:p15-measurable-map} and then \(\mathbb P\); it does not presuppose
that \(h\) is invertible. It is the law \(\mathbb P_h\) of
\cref{def:p15-pushforward-law}, by definition.

The formula has no Jacobian factor, because the measure itself is
transported rather than compared with Lebesgue measure. The calculus rule
that substitutes \(y=\varphi(x)\) inside a Lebesgue integral and inserts
\(|\det D\varphi|\) is a different theorem, stated in this part for a
\(C^1\) diffeomorphism in \cref{thm:p15-w02b-p15-042-diffeomorphism} and
for affine maps in \cref{thm:p15-affine-lebesgue}.
\index{change of variables formula}\index{pushforward measure!change of variables}
\end{hbconvention}

\begin{hbcorollary}{thm:p15-w02c-p15-001-change-of-variables}{proved=annex:p15-w02c-p15-001-change-of-variables}{Change of variables formula}
\par\nobreak
Let \(X\) be a random element of \((S,\mathcal S)\) with distribution
\(\mu\), that is, \(\mu(A)=\mathbb P(X\in A)\) for \(A\in\mathcal S\).
If \(f\) is a measurable function from \((S,\mathcal S)\) to
\((\mathbb R,\mathcal B(\mathbb R))\) such that \(f\geq0\) or
\(\mathbb E|f(X)|<\infty\), then
\[
 \mathbb E f(X)=\int_S f(y)\,\mu(dy).
\]
Equivalently, writing \(h\) for \(X\),
\(\int_\Omega f(h(\omega))\,d\mathbb P(\omega)
=\int_S f(y)\,d(\mathbb P\circ h^{-1})(y)\), and \(\mathbb E|f(X)|<\infty\)
holds exactly when \(\int_S|f(y)|\,\mu(dy)<\infty\).
\end{hbcorollary}

Durrett draws the practical consequence: expected values of functions of
real random variables can be computed as integrals on the real line, where
calculus is available \cite[sec.~1.6.3, pp.~32--33]{DurrettPTEAuthor2019}.
The next result does this for the uniform variable.

\begin{hbcorollary}{thm:p15-w02c-p15-001-uniform-integrals}{proved=annex:p15-w02c-p15-001-uniform-integrals}{Expectations of the uniform variable as interval integrals}
\par\nobreak
On the uniform probability space \((\Omega_U,\mathcal F_U,\mathbb U)\) of
\cref{thm:p15-lebesgue-unit-interval}, the law of the identity variable
\(U\) is
\[
 \mathbb U_U(B)=\lambda(B\cap(0,1))=\int_B\mathbf1_{(0,1)}\,d\lambda
 \qquad(B\in\mathcal B(\mathbb R)),
\]
so \(\mathbf1_{(0,1)}\) is a density for it. Let \(h:\mathbb R\to\mathbb R\)
be Borel with \(h\geq0\) on \([0,1]\), or with
\(\int_0^1|h(x)|\,dx<\infty\). Then
\[
 \mathbb E[h(U)]=\int_0^1h(x)\,dx,\qquad
 \mathbb E[|h(U)|]=\int_0^1|h(x)|\,dx,
\]
the first equality holding in \([0,\infty]\) in the nonnegative case and
between finite numbers in the integrable case.
\end{hbcorollary}

\begin{hbexample}{def:p15-w02c-p15-001-uniform-moments-as-integrals}{Two interval integrals computed from the uniform moments}
Take \(h(x)=x\) and \(h(x)=x^2\). Both are continuous, hence Borel by
\cref{thm:p15-measurable-operations}, and both satisfy \(|h|\leq1\) on
\([0,1]\), so \(\int_0^1|h(x)|\,dx\leq1\) by part~(3) of
\cref{thm:p15-w02c-p15-001-interval-integrals} and monotonicity.
\Cref{thm:p15-w02b-p15-001-uniform-moments} gives \(\mathbb E[U]=1/2\)
and \(\mathbb E[U^2]=1/3\), computed there from simple lower and upper
bounds on the uniform space. By
\cref{thm:p15-w02c-p15-001-uniform-integrals},
\[
 \int_0^1x\,dx=\frac12,\qquad\int_0^1x^2\,dx=\frac13,\qquad
 \int_0^1\left(x-\frac12\right)^2dx=\operatorname{Var}(U)=\frac1{12}.
\]
The last identity applies the corollary to the nonnegative function
\((x-1/2)^2\). No derivative or antiderivative enters: the
values come from the integral of
\cref{def:p15-nonnegative-integral} on the uniform space, transferred to
the line by the change of variables formula.

The same transfer works in two dimensions. By
\cref{thm:p15-finite-product-measures}, the indicator of the unit square
\((0,1]^2\) has \(\int\mathbf1_{(0,1]^2}(x)\,dx=\lambda^2((0,1]^2)=1\),
where \(x=(x_1,x_2)\) is the variable of integration in
case~(a) of \cref{def:p15-w02c-p15-001-dx-notation}.
\end{hbexample}

\begin{hbwarning}{def:p15-w02c-p15-001-change-of-variables-limits}{Limits of the change of variables formula}
\par\nobreak
\begin{enumerate}
\item \emph{A signed function with two infinite parts.} The formula is
      stated for \(f\geq0\) or \(\mathbb E|f(X)|<\infty\). If both
      \(\mathbb E[f(X)^+]\) and \(\mathbb E[f(X)^-]\) are infinite,
      neither side is defined, and the formula asserts nothing. The
      variable \(X(s,n)=s2^n\) in the second part of
      \cref{def:p15-computed-expectations}, with \(f(y)=y\), is an
      instance.
\item \emph{Knowledge of the law, not of the sample space.} Both sides
      depend on \(X\) only through \(\mathbb P_X\). Two random elements
      with the same law, possibly on different probability spaces, have
      the same expectations \(\mathbb E f(X)\); the formula cannot
      distinguish them, and an expectation that depends on more than the
      law of \(X\), such as \(\mathbb E[f(X)Y]\) for a second variable
      \(Y\), needs the joint law of \((X,Y)\).
\item \emph{A density is a further step.} The value-space integral
      \(\int_S f\,d\mathbb P_X\) becomes an ordinary \(dx\) integral only
      when \(\mathbb P_X\) has a density relative to \(\lambda^d\), by
      \cref{thm:p15-density-integrals}; a law with atoms, such as the
      Dirac law \(\delta_c\), has none.
\end{enumerate}
\end{hbwarning}

\section{The four-step method}
\label{sec:p15-w02c-p15-001-four-step}

Durrett proves the change of variables formula in four cases of increasing
generality, which follow the way the integral was built, and asks the
reader to note the method, since it is used several times later
\cite[proof of thm.~1.6.9, p.~32]{DurrettPTEAuthor2019}. The method is
stated here as a procedure and as a theorem that carries out its last
three steps once and for all.

\begin{hbmethod}{def:p15-w02c-p15-001-four-step-method}{Proving an integral identity in four steps}
\emph{Purpose.} To show that two expressions \(\Phi(f)\) and \(\Psi(f)\),
each built from integrals of a function \(f\) on a measurable space
\((S,\mathcal S)\), agree for every nonnegative measurable \(f\) and for
every \(f\) for which they are finite.

\emph{Input.} The two expressions, defined for measurable
\(f:S\to[0,\infty]\), with values in \([0,\infty]\).

\emph{Procedure.}
\begin{enumerate}
\item \emph{Indicators.} Verify \(\Phi(\mathbf1_B)=\Psi(\mathbf1_B)\) for
      every \(B\in\mathcal S\). This is the only step that uses the
      particular expressions; it is usually an unwinding of definitions.
\item \emph{Simple functions.} Check that each expression is additive and
      commutes with multiplication by finite constants \(c\geq0\); then the
      identity passes to nonnegative simple functions
      \(\sum_ja_j\mathbf1_{B_j}\).
\item \emph{Nonnegative functions.} Check that each expression is
      continuous along increasing sequences, usually by the monotone
      convergence theorem \cref{thm:p15-monotone-integral}; then the
      identity passes to every nonnegative measurable \(f\) through simple
      functions increasing to \(f\).
\item \emph{Integrable functions.} For a real measurable \(f\), apply
      step~3 to \(|f|\), \(f^+\) and \(f^-\). The finiteness conditions
      \(\Phi(|f|)<\infty\) and \(\Psi(|f|)<\infty\) then coincide, and when
      they hold the two signed values \(\Phi(f^+)-\Phi(f^-)\) and
      \(\Psi(f^+)-\Psi(f^-)\) agree.
\end{enumerate}

\emph{Output.} The identity for all nonnegative measurable functions,
in \([0,\infty]\), and for all functions satisfying the finiteness
condition of step~4, as real numbers.

\emph{Side conditions.} Step~1 must cover all of \(\mathcal S\). If it
is checked only on a generating class, a uniqueness theorem has to fill
the gap first, and it needs its hypotheses:
\cref{thm:p15-w02c-p15-002-sigma-finite-uniqueness} asks for a pi-system
with a finite exhaustion, and \cref{def:p15-w02c-p15-002-uniqueness-failure}
shows two measures that agree on a generating pi-system but differ. Step~4
gives nothing for a function with \(\Phi(f^+)=\Phi(f^-)=\infty\); see the
first item of \cref{def:p15-w02c-p15-001-change-of-variables-limits}.

\emph{Source.} Durrett defines the integral in four steps, namely for
simple, bounded, nonnegative and general functions, and remarks that the
same sequence of steps helps in proving integration formulas, pointing to
his proofs of Theorems~1.6.9 and~1.7.2
\cite[sec.~1.4, p.~18]{DurrettPTEAuthor2019}. Those proofs run instead
through indicators, simple functions, nonnegative functions and
integrable functions: for the change of variables formula
\cite[pp.~32--33]{DurrettPTEAuthor2019}, and for Fubini's theorem, his
Theorem~1.7.2, stated on page~38 and proved on page~39
\cite[thm.~1.7.2, pp.~38--39]{DurrettPTEAuthor2019}.
The procedure above is the proof scheme. The integral of this handbook is
built in three stages (\cref{def:p15-simple-integral},
\cref{def:p15-nonnegative-integral}, \cref{def:p15-integrable-functions}),
and \cref{thm:p15-integral-source-bridge} identifies it with Durrett's
bounded and nonnegative stages.
\index{four-step method}\index{integral!four-step method}
\end{hbmethod}

\begin{hbtheorem}{thm:p15-w02c-p15-001-four-step}{proved=annex:p15-w02c-p15-001-four-step}{The four-step method as a theorem}
\par\nobreak
Let \((S,\mathcal S)\) be a measurable space and let \(\mathcal M_+\) be
the set of measurable functions \(S\to[0,\infty]\). Let
\(\Phi,\Psi:\mathcal M_+\to[0,\infty]\), and suppose that each
\(\Theta\in\{\Phi,\Psi\}\) satisfies:
\begin{enumerate}
\item[(a)] \(\Theta(f+g)=\Theta(f)+\Theta(g)\) and
      \(\Theta(cf)=c\,\Theta(f)\) for \(f,g\in\mathcal M_+\) and finite
      \(c\geq0\), with \(0\cdot\infty=0\);
\item[(b)] if \(f_n\in\mathcal M_+\), \(f_n\leq f_{n+1}\) and
      \(f_n\to f\) pointwise, then \(\Theta(f_n)\to\Theta(f)\).
\end{enumerate}
If \(\Phi(\mathbf1_B)=\Psi(\mathbf1_B)\) for every \(B\in\mathcal S\),
then:
\begin{enumerate}
\item \(\Phi(f)=\Psi(f)\) for every \(f\in\mathcal M_+\);
\item for a real measurable \(f\) on \(S\), \(\Phi(|f|)<\infty\) exactly
      when \(\Psi(|f|)<\infty\), and then \(\Phi(f^+)\), \(\Phi(f^-)\),
      \(\Psi(f^+)\), \(\Psi(f^-)\) are finite and
      \(\Phi(f^+)-\Phi(f^-)=\Psi(f^+)-\Psi(f^-)\).
\end{enumerate}
\end{hbtheorem}

For a random element \(X\) of \((S,\mathcal S)\), the pair
\(\Phi(h)=\mathbb E[h(X)]\) and \(\Psi(h)=\int_Sh\,d\mathbb P_X\)
satisfies the hypotheses, and the theorem yields the change of variables
formula; the annex proof of \cref{thm:p15-w02c-p15-001-change-of-variables}
records this derivation. Other instances on this part's route are the
density formula \(\int h\,d\eta=\int hf\,d\lambda^d\) of
\cref{thm:p15-density-integrals}, with \(\Phi(h)=\int h\,d\eta\) and
\(\Psi(h)=\int hf\,d\lambda^d\), and the iterated integrals of Fubini's
theorem \cref{thm:p15-fubini}, where step~1 is the product-measure
identity for a measurable set.

\section{A direct proof of Cauchy--Schwarz}
\label{sec:p15-w02c-p15-001-cauchy-schwarz}

\Cref{thm:p15-holder-integral} obtains the Cauchy--Schwarz inequality as
the case \(p=q=2\) of H\"older's inequality, which Durrett states as his
Theorem~1.5.2 and proves through the elementary bound
\(xy\leq x^p/p+y^q/q\) \cite[thm.~1.5.2, p.~25]{DurrettPTEAuthor2019}.
In the remark that follows, Durrett names the case \(p=q=2\) and gives a
direct argument for it \cite[remark after thm.~1.5.2, p.~26]{DurrettPTEAuthor2019}:
for every real \(\theta\),
\[
 0\leq\int(f+\theta g)^2\,d\mu
 =\int f^2\,d\mu+\theta\,2\int fg\,d\mu+\theta^2\int g^2\,d\mu ,
\]
so the quadratic \(a\theta^2+b\theta+c\) on the right has at most one real
root, and the formula for the roots of a quadratic then forces
\(b^2-4ac\leq0\). The theorem below states what that argument proves, with
the hypotheses that make each step meaningful; the annex proof writes out
the step from nonnegativity to the discriminant and the case \(a=0\), in
which the right-hand side is not a quadratic.

\begin{hbtheorem}{thm:p15-w02c-p15-001-cauchy-schwarz}{proved=annex:p15-w02c-p15-001-cauchy-schwarz}{Cauchy--Schwarz by the discriminant}
\par\nobreak
Let \((\Omega,\mathcal F,\mu)\) be a measure space and let
\(f,g:\Omega\to\mathbb R\) be measurable with
\(\int f^2\,d\mu<\infty\) and \(\int g^2\,d\mu<\infty\). Put
\[
 a\defeq\int g^2\,d\mu,\qquad b\defeq2\int fg\,d\mu,\qquad
 c\defeq\int f^2\,d\mu .
\]
Then \(fg\) is integrable, so \(a\), \(b\) and \(c\) are real numbers, and:
\begin{enumerate}
\item for every real \(\theta\),
      \(\int(f+\theta g)^2\,d\mu=a\theta^2+b\theta+c\geq0\);
\item \(b^2-4ac\leq0\), that is,
      \[
       \left(\int fg\,d\mu\right)^2\leq\int f^2\,d\mu\int g^2\,d\mu ;
      \]
\item \(\int|fg|\,d\mu\leq\|f\|_2\|g\|_2\), with the norms of
      \cref{def:p15-lp-convergence}.
\end{enumerate}
\end{hbtheorem}

\begin{hbexample}{def:p15-w02c-p15-001-cauchy-schwarz-uniform}{The first two uniform moments}
On the uniform probability space of \cref{thm:p15-lebesgue-unit-interval},
take \(f\equiv1\) and \(g=U\). Both are square integrable, since
\(\mathbb E[1]=1\) and \(\mathbb E[U^2]=1/3\) by
\cref{thm:p15-w02b-p15-001-uniform-moments}. Here \(a=1/3\),
\(b=2\,\mathbb E[U]=1\) and \(c=1\), so \(b^2-4ac=1-4/3=-1/3\leq0\), and
part~(2) reads \((\mathbb E[U])^2=1/4\leq1/3=\mathbb E[1]\,\mathbb E[U^2]\).
The quadratic of part~(1) is
\(\mathbb E[(1+\theta U)^2]=\theta^2/3+\theta+1\). Its value at
\(\theta=-3/2\) is \(3/4-3/2+1=1/4\), and
\(\theta^2/3+\theta+1=(\theta+3/2)^2/3+1/4\geq1/4\) for every real
\(\theta\), so it has no real root.
\end{hbexample}

\begin{hbwarning}{def:p15-w02c-p15-001-cauchy-schwarz-limits}{Where the quadratic argument needs care}
\par\nobreak
\begin{enumerate}
\item \emph{A degenerate quadratic.} If \(\int g^2\,d\mu=0\), so that
      \(g=0\) a.e., then \(a=0\) and part~(1) is a statement about the
      affine function \(b\theta+c\); the formula for the roots, which
      divides by \(2a\), does not apply. The conclusion survives because a
      nonnegative affine function of \(\theta\) on the whole line must be
      constant, so \(b=0\). The annex proof treats this case separately.
\item \emph{Infinite second moments.} The expansion of
      \(\int(f+\theta g)^2\,d\mu\) needs \(a\), \(b\) and \(c\) to be real.
      On the uniform space, let \(f=\sqrt{q(U)}\), that is
      \(f(u)=u^{-1/2}\), with \(q\) the reciprocal function of
      \cref{thm:p15-w02c-p15-001-reciprocal-density}, and let \(g=-f\).
      By \cref{thm:p15-w02c-p15-001-uniform-integrals} and that theorem,
      \(\mathbb E[f^2]=\int_0^1q(x)\,dx=\infty\), and
      \(\int fg\,d\mathbb U=-\mathbb E[f^2]=-\infty\). At \(\theta=1\) the
      left side is \(\int0\,d\mathbb U=0\), while the right side would be
      \(\infty-\infty+\infty\), which has no value. H\"older's inequality
      in \cref{thm:p15-holder-integral} still holds in such cases, with
      its product convention; the direct argument does not reach them.
\end{enumerate}
\end{hbwarning}

\section{Chebyshev's inequality: the names and the squared form}
\label{sec:p15-w02c-p15-001-chebyshev}

\begin{hbconvention}{def:p15-w02c-p15-001-chebyshev-names}{The names of Chebyshev's inequality}
Durrett's Theorem~1.6.4, titled \emph{Chebyshev's inequality}, is the
general bound: if \(\varphi:\mathbb R\to\mathbb R\) has \(\varphi\geq0\),
\(A\subseteq\mathbb R\) and \(i_A=\inf\{\varphi(y):y\in A\}\), then
\[
 i_A\,\mathbb P(X\in A)\leq\mathbb E(\varphi(X);X\in A)\leq\mathbb E\varphi(X)
\]
\cite[thm.~1.6.4, p.~29]{DurrettPTEAuthor2019}. It is the first display of
\cref{thm:p15-expectation-bounds}, where \(\varphi\) is Borel, \(A\) is a
nonempty Borel set and the infimum is written \(c_A\). In the remark after
its proof, Durrett reports a second usage: some authors call the general
bound \emph{Markov's inequality} and keep the name Chebyshev's inequality
for the special case \(\varphi(x)=x^2\), \(A=\{x:|x|\geq a\}\), which is
his display~(1.6.1),
\[
 a^2\,\mathbb P(|X|\geq a)\leq\mathbb E X^2
\]
\cite[remark after thm.~1.6.4, p.~30]{DurrettPTEAuthor2019}.
So one inequality and its special cases carry the two names in different
ways. Under Durrett's title, the general bound is Chebyshev's. Under the
usage he reports, the general bound is Markov's, and Chebyshev's
is~(1.6.1). In \cref{thm:p15-expectation-bounds}, the name ``Markov's
bound'' is given to a different special case,
\(\mathbb P(Y\geq r)\leq\mathbb E[Y]/r\) for nonnegative \(Y\) and
\(r>0\), which is the general bound with \(\varphi(y)=\max(y,0)\) and
\(A=[r,\infty)\), and the name Chebyshev is given to the centred variance
form. The mathematics does not depend on the choice of names; a citation
of ``Chebyshev's inequality'' or ``Markov's inequality'' has to be read
against the statement it names.
\index{Chebyshev's inequality!names}\index{Markov's inequality!names}
\end{hbconvention}

\begin{hbcorollary}{thm:p15-w02c-p15-001-chebyshev-square}{proved=annex:p15-w02c-p15-001-chebyshev-square}{The squared form~(1.6.1) and its centred version}
\par\nobreak
Let \(X\) be a real random variable on \((\Omega,\mathcal F,\mathbb P)\).
\begin{enumerate}
\item For every real \(a\geq0\),
      \(a^2\,\mathbb P(|X|\geq a)\leq\mathbb E X^2\), in \([0,\infty]\).
\item If \(\mathbb E X^2<\infty\), then for every real \(r>0\),
      \(\mathbb P(|X-\mathbb E X|\geq r)\leq\operatorname{Var}(X)/r^2\); this
      is part~(1) applied to the random variable \(X-\mathbb E X\).
\item For \(a<0\) the inequality of part~(1) can fail: for the constant
      variable \(X\equiv0\), the left side is \(a^2>0\) and the right side
      is \(0\).
\end{enumerate}
\end{hbcorollary}

Durrett's remark does not state the range of \(a\) in~(1.6.1). The
derivation from his Theorem~1.6.4 needs \(i_A=a^2\) for
\(A=\{x:|x|\geq a\}\), and that identity holds exactly when \(a\geq0\):
for \(a<0\) the set \(A\) is the whole line and \(i_A=0\). Part~(3) shows
that the range \(a\geq0\) cannot be dropped.

\begin{hbexample}{def:p15-w02c-p15-001-chebyshev-uniform}{The bound for a uniform variable}
Let \(U\) be the uniform variable, with \(\mathbb E[U]=1/2\) and
\(\operatorname{Var}(U)=1/12\)
(\cref{def:p15-w02c-p15-001-uniform-moments-as-integrals}). Take
\(r=2/5\). The event \(\{|U-1/2|\geq2/5\}\) is
\(\{U\leq1/10\}\cup\{U\geq9/10\}\), and the law of \(U\) in
\cref{thm:p15-w02c-p15-001-uniform-integrals} gives it probability
\(\lambda((0,1/10]\cup[9/10,1))=1/10+1/10=1/5\). Part~(2) of
\cref{thm:p15-w02c-p15-001-chebyshev-square} bounds this probability by
\((1/12)/(4/25)=25/48\). The bound holds and is far from attained; for
\(r\leq1/\sqrt{12}\) it is at least one and says nothing.
\end{hbexample}

\section{Interval integrals of affine functions and of the reciprocal}
\label{sec:p15-w02c-p15-001-affine-integrals}

Several examples in this part quote values of one-variable integrals,
such as \(\int_0^1(2x-1)\,dx=0\), which elementary calculus obtains from an
antiderivative. The integrals here are Lebesgue integrals in the sense of
case~(b) of \cref{def:p15-w02c-p15-001-dx-notation}, and no fundamental
theorem of calculus is proved in this part. The next results compute the
values that the examples of \cref{thm:p15-radon-nikodym-examples} use,
directly from the integral of \cref{def:p15-nonnegative-integral} and the
signed integral of \cref{def:p15-integrable-functions}: an affine function
is squeezed between step functions, and the reciprocal function is
bounded below on dyadic pieces of \((0,a)\).

\begin{hblemma}{thm:p15-w02c-p15-001-affine-integral}{proved=annex:p15-w02c-p15-001-affine-integral}{The integral of an affine function over an interval}
\par\nobreak
Let \(a\leq b\) and \(\alpha,\beta\) be real. Then
\(\int_a^b|\alpha x+\beta|\,dx<\infty\), and
\[
 \int_a^bx\,dx=\frac{b^2-a^2}2,\qquad
 \int_a^b(\alpha x+\beta)\,dx=\alpha\,\frac{b^2-a^2}2+\beta(b-a).
\]
\end{hblemma}

\begin{hbcorollary}{thm:p15-w02c-p15-001-signed-density-values}{proved=annex:p15-w02c-p15-001-signed-density-values}{The integrals of the signed density \(2x-1\)}
\par\nobreak
Let \(h(x)=2x-1\). Then
\[
 \int_0^1(2x-1)\,dx=0,\qquad
 \int_{1/2}^1(2x-1)\,dx=\frac14,\qquad
 \int_0^{1/2}(1-2x)\,dx=\frac14 ,
\]
and
\(\int_{[0,1]}h^+\,d\lambda=\int_{[0,1]}h^-\,d\lambda=1/4\) and
\(\int_{[0,1]}|h|\,d\lambda=1/2\). These are the values used in part~(2)
of \cref{thm:p15-radon-nikodym-examples}.
\end{hbcorollary}

\begin{hbtheorem}{thm:p15-w02c-p15-001-reciprocal-density}{proved=annex:p15-w02c-p15-001-reciprocal-density}{The reciprocal density near zero}
\par\nobreak
Let \(q:\mathbb R\to[0,\infty)\) be given by \(q(x)=1/x\) for
\(0<x<1\) and \(q(x)=0\) otherwise. Then \(q\) is Borel, and:
\begin{enumerate}
\item for every integer \(n\geq2\),
      \(\int_{[1/n,1)}q\,d\lambda\leq n-1<\infty\);
\item for every real \(a\in(0,1]\),
      \(\int_{(0,a)}q\,d\lambda=\infty\); in particular
      \(\int_0^1q(x)\,dx=\infty\).
\end{enumerate}
These are the two computations behind part~(4) of
\cref{thm:p15-radon-nikodym-examples}: the sets \([1/n,1)\), \(n\geq2\),
cover \((0,1)\) with finite mass, and every interval \((0,a)\) has
infinite mass although its Lebesgue measure \(a\) is as small as one
likes.
\end{hbtheorem}

The bound in part~(1) is an upper bound, which is all that the covering
argument for sigma-finiteness requires. Part~(2) is a lower bound by a
divergent sum, so no value of an antiderivative is needed. The function
\(q\) is also the limitation of the step-function method of
\cref{thm:p15-w02c-p15-001-affine-integral}: it is unbounded on \((0,1)\),
so no step function with finitely many steps lies above it there.

\section{Lower absolute moments and the Bernoulli law}
\label{sec:p15-w02c-p15-001-lower-moments}

Two small facts about moments recur in later chapters: a finite absolute
moment of some order forces finite absolute moments of every smaller
positive order, and a variable that takes the values \(1\) and \(0\) has
mean \(p\) and variance \(p(1-p)\). Durrett sets the first as an exercise
\cite[exercise~1.6.11, p.~37]{DurrettPTEAuthor2019} and records the
second as his Example~1.6.12
\cite[example~1.6.12, p.~34]{DurrettPTEAuthor2019}. Both are proved here
from the integral of \cref{sec:p15-integration-expectation}, and the
limits of the first are shown on two explicit measures.

\emph{Inputs.} A real random variable \(X\) on a probability space
\((\Omega,\mathcal F,\mathbb P)\); the expectation and variance of
\cref{def:p15-expectation-integral}; H\"older's inequality
(\cref{thm:p15-holder-integral}); the a.e.\ rules of
\cref{thm:p15-nonnegative-integral-laws}; the variance identity of
\cref{thm:p15-expectation-bounds}. \emph{Conclusions.} For real
\(0<j<k\), \(\mathbb E|X|^j\leq(\mathbb E|X|^k)^{j/k}\) in
\([0,\infty]\), so a finite \(k\)-th absolute moment gives a finite
\(j\)-th one; the bound fails without unit total mass, and the converse
fails; a Bernoulli variable with parameter \(p\) has
\(\mathbb EX=\mathbb EX^2=p\) and \(\operatorname{Var}(X)=p(1-p)\).

\begin{hbtheorem}{thm:p15-w02c-p15-001-lower-moments}{proved=annex:p15-w02c-p15-001-lower-moments}{Finite higher absolute moments give finite lower ones}
\par\nobreak
Let \(X\) be a real random variable on \((\Omega,\mathcal F,\mathbb P)\)
and let \(j,k\) be real numbers with \(0<j<k\). Then \(|X|^j\) and
\(|X|^k\) are nonnegative random variables, and, with
\(\infty^{j/k}\defeq\infty\),
\[
 \mathbb E|X|^j\leq\left(\mathbb E|X|^k\right)^{j/k}
 \qquad\text{in }[0,\infty].
\]
In particular, if \(\mathbb E|X|^k<\infty\), then \(\mathbb E|X|^j<\infty\).
Equivalently, with \(\infty^{1/r}\defeq\infty\) for \(r>0\),
\[
 \left(\mathbb E|X|^j\right)^{1/j}\leq\left(\mathbb E|X|^k\right)^{1/k},
\]
that is, \(\|X\|_j\leq\|X\|_k\) in the notation of
\cref{def:p15-lp-convergence} whenever \(1\leq j<k\).
\end{hbtheorem}

Durrett's exercise states the finiteness claim and the displayed bound
for \(0<j<k\) under the hypothesis \(\mathbb E|X|^k<\infty\)
\cite[exercise~1.6.11, p.~37]{DurrettPTEAuthor2019}; it does not say
whether \(j\) and \(k\) are integers. The theorem covers all real orders,
the integer orders among them. When \(\mathbb E|X|^k=\infty\) the bound
holds trivially.

\begin{hbtheorem}{thm:p15-w02c-p15-001-lower-moments-limits}{proved=annex:p15-w02c-p15-001-lower-moments-limits}{What the moment bound needs, and what it does not give}
\par\nobreak
\begin{enumerate}
\item \emph{Unit total mass is needed.} Let \(\#\) be counting measure on
      all subsets of \(I=\{1,2,\ldots\}\) and \(f(i)=1/i\). Then
      \(\int f^2\,d\#\leq2\) while \(\int f\,d\#=\infty\). So on a
      space of infinite measure a finite integral of \(|f|^2\) does not give
      a finite integral of \(|f|\).
\item \emph{The converse fails.} On \(I\) with all subsets, let
      \(\mathbb P(A)\defeq\sum_{n\in A}2^{-n}\). This is a probability
      measure, and the variable \(X(n)=2^n\) has
      \[
       \mathbb E X^{1/2}=\frac1{\sqrt2-1}<\infty,\qquad
       \mathbb EX=\infty .
      \]
      So a finite moment of a smaller order does not give a finite moment
      of a larger one.
\item \emph{The bound is attained.} If \(|X|=c\) a.s.\ for a real
      constant \(c\geq0\), then
      \(\mathbb E|X|^j=c^j=(\mathbb E|X|^k)^{j/k}\).
\end{enumerate}
\end{hbtheorem}

\begin{hbdefinition}{def:p15-w02c-p15-001-bernoulli}{The Bernoulli law}
Let \(0\leq p\leq1\). A real random variable \(X\) has the
\emph{Bernoulli distribution with parameter \(p\)} when
\(\mathbb P(X=1)=p\) and \(\mathbb P(X=0)=1-p\)
\cite[example~1.6.12, p.~34]{DurrettPTEAuthor2019}. For every event
\(A\in\mathcal F\), the indicator \(\mathbf1_A\) has this law with
\(p=\mathbb P(A)\), since \(\{\mathbf1_A=1\}=A\) and
\(\{\mathbf1_A=0\}=A^c\).
\index{Bernoulli distribution}
\end{hbdefinition}

\begin{hbtheorem}{thm:p15-w02c-p15-001-bernoulli-moments}{proved=annex:p15-w02c-p15-001-bernoulli-moments}{Mean and variance of a Bernoulli variable}
\par\nobreak
Let \(X\) have the Bernoulli distribution with parameter \(p\). Then
\(X\) is integrable, \(X^2=X\) almost surely, and
\[
 \mathbb EX=p,\qquad \mathbb EX^2=p,\qquad
 \operatorname{Var}(X)=\mathbb EX^2-(\mathbb EX)^2=p(1-p).
\]
Moreover \(p(1-p)=\tfrac14-(p-\tfrac12)^2\leq\tfrac14\), with equality
exactly when \(p=1/2\), and \(\operatorname{Var}(X)=0\) exactly when
\(p\in\{0,1\}\).
\end{hbtheorem}

Durrett computes the mean as \(p\cdot1+(1-p)\cdot0\) and uses the identity
\(X^2=X\) \cite[example~1.6.12, p.~34]{DurrettPTEAuthor2019}. On a
general probability space the hypotheses fix the values of \(X\) only off
the null event \(\{X\notin\{0,1\}\}\), so \(X^2=X\) is an almost-sure
identity, and the computation passes through the a.e.\ rules of
\cref{thm:p15-nonnegative-integral-laws}.

\begin{hbexample}{def:p15-w02c-p15-001-lower-moments-example}{An indicator of the uniform variable}
Let \(U\) be the uniform variable of
\cref{thm:p15-w02c-p15-001-uniform-integrals}, with
\(\mathbb E[U]=1/2\) and \(\mathbb E[U^2]=1/3\)
(\cref{def:p15-w02c-p15-001-uniform-moments-as-integrals}).
\begin{enumerate}
\item The event \(A=\{U\leq1/3\}\) has
      \(\mathbb P(A)=\lambda((0,1/3])=1/3\) by the law of \(U\) in
      \cref{thm:p15-w02c-p15-001-uniform-integrals}, so \(\mathbf1_A\) is
      Bernoulli with parameter \(1/3\). By
      \cref{thm:p15-w02c-p15-001-bernoulli-moments}, its mean is \(1/3\)
      and its variance is \((1/3)(2/3)=2/9\).
\item With \(j=1\) and \(k=2\),
      \cref{thm:p15-w02c-p15-001-lower-moments} for \(U\) says
      \(1/2\leq(1/3)^{1/2}\), which holds since \(1/4\leq1/3\). For
      \(\mathbf1_A\) it says \(1/3\leq(1/3)^{1/2}\), which holds strictly
      since \(1/9<1/3\). Neither variable has an almost surely constant
      absolute value, and in both cases the bound is not attained.
\end{enumerate}
\end{hbexample}

\begin{hbwarning}{def:p15-w02c-p15-001-lower-moments-warning}{Where the moment bound stops}
The comparison of moments uses the unit total mass of \(\mathbb P\)
through the factor \(\|1\|_q=1\) of H\"older's inequality in its proof.
Part~(1) of \cref{thm:p15-w02c-p15-001-lower-moments-limits} shows that
the conclusion fails for counting measure on an infinite set. Part~(2)
shows that the comparison runs only from higher to lower orders. For the
Bernoulli law, the value \(p(1-p)\) belongs to the values \(0\) and \(1\);
a variable that takes two other values has a different variance, computed
from the general identity of \cref{thm:p15-expectation-bounds}.
\end{hbwarning}

The moment comparison of \cref{thm:p15-w02c-p15-001-lower-moments} is the
case \(g=1\) of H\"older's inequality (\cref{thm:p15-holder-integral});
with \(j=1\) and \(k=2\) it is also the inequality
\((\mathbb E|X|)^2\leq\mathbb E(X^2)\), which Jensen's inequality
(\cref{thm:p15-jensen-integral}) with \(\varphi(x)=x^2\) gives when
\(\mathbb E(X^2)<\infty\). The Bernoulli law is the law of the indicator of a
single event; \cref{sec:p15-finite-probability} treats such indicators on
finite outcome spaces.

\section{The hypotheses of Jensen's inequality}
\label{sec:p15-w02c-p15-001-jensen-hypotheses}

Durrett states Jensen's inequality twice. For integrals, his
Theorem~1.5.1 assumes that \(\mu\) is a probability measure and that
\(f\) and \(\varphi(f)\) are integrable
\cite[thm.~1.5.1, p.~25]{DurrettPTEAuthor2019}. For expected values,
his Theorem~1.6.2 asserts \(\mathbb E(\varphi(X))\geq\varphi(\mathbb EX)\)
provided both expectations exist, which he spells out as
\(\mathbb E|X|<\infty\) and \(\mathbb E|\varphi(X)|<\infty\)
\cite[thm.~1.6.2, p.~29]{DurrettPTEAuthor2019}. The handbook statement
\cref{thm:p15-jensen-integral} carries the same three hypotheses. This
section shows what each of them does. The total mass one cannot be dropped
on any finite measure space; on a space of infinite total mass the
conclusion may fail or hold, so the total mass alone does not decide it.
Integrability of \(\varphi(X)\) can be dropped once \(\mathbb E\varphi(X)\)
is read as an extended integral, because that integral then always
exists and the only case the hypothesis excludes is the value \(\infty\).
Integrability of \(X\) cannot be dropped, since without it the left side
\(\varphi(\mathbb EX)\) may have no meaning.

\emph{Inputs.} The integral and its existence conventions
(\cref{def:p15-nonnegative-integral,def:p15-integrable-functions});
convex functions (\cref{def:p15-convex-integral}) and the supporting
affine function constructed in the proof of
\cref{thm:p15-jensen-integral}; countable atomic models
(\cref{thm:p15-countable-atomic}); Lebesgue measure \(\lambda\)
(\cref{thm:p15-lebesgue-unit-interval}). \emph{Conclusions.} On a
measure space of finite total mass, Jensen's conclusion holds for every
admissible pair \((f,\varphi)\) exactly when the total mass is one; two
measures of infinite total mass, on one of which it fails and on the
other of which it holds; the extended form for a convex function of an
integrable variable; and explicit variables showing that \(\varphi(X)\)
may fail to be integrable when \(X\) is, and that \(\mathbb EX\) may be
undefined.

Durrett's own illustration of the direction of the inequality, a
two-point law, and his Figure~1.6 are worked in
\cref{def:p15-w02c-p15-005-jensen-direction}
\cite[pp.~29--30]{DurrettPTEAuthor2019}; the special cases
\(|\mathbb EX|\leq\mathbb E|X|\) and
\((\mathbb EX)^2\leq\mathbb E(X^2)\) that he names after
Theorem~1.6.2 are \cref{thm:p15-w02c-p15-005-jensen-special}.

\begin{hbtheorem}{thm:p15-w02c-p15-001-jensen-normalization}{proved=annex:p15-w02c-p15-001-jensen-norm}{Total mass one is necessary on a finite measure space}
\par\nobreak
Let \((\Omega,\mathcal F,\mu)\) be a measure space with
\(\mu(\Omega)<\infty\). The following are equivalent.
\begin{enumerate}
\item \(\mu(\Omega)=1\).
\item For every real integrable \(f\) and every convex
      \(\varphi:\mathbb R\to\mathbb R\) such that \(\varphi(f)\) is
      integrable,
      \(\varphi\bigl(\int f\,d\mu\bigr)\leq\int\varphi(f)\,d\mu\).
\end{enumerate}
More precisely, with \(f\equiv0\): if \(\mu(\Omega)<1\), the constant
\(\varphi\equiv1\) violates the inequality in~(2), and if
\(\mu(\Omega)>1\), the constant \(\varphi\equiv-1\) violates it.
\end{hbtheorem}

\begin{hbtheorem}{thm:p15-w02c-p15-001-jensen-infinite-mass}{proved=annex:p15-w02c-p15-001-jensen-infinite-mass}{Infinite total mass does not decide Jensen's inequality}
\par\nobreak
Let \((\Omega,\mathcal F,\mu)\) be a measure space with
\(\mu(\Omega)=\infty\).
\begin{enumerate}
\item If some \(A\in\mathcal F\) has \(1<\mu(A)<\infty\), then
      \(f=\mathbf1_A\) and \(\varphi(x)=x^2\) satisfy the hypotheses
      in~(2) of \cref{thm:p15-w02c-p15-001-jensen-normalization} and
      \[
       \varphi\Bigl(\int f\,d\mu\Bigr)=\mu(A)^2>\mu(A)=\int\varphi(f)\,d\mu .
      \]
      On \((\mathbb R,\mathcal B(\mathbb R),\lambda)\) the set
      \(A=(0,2]\) gives \(4>2\).
\item If \(\mu(A)\in\{0,\infty\}\) for every \(A\in\mathcal F\), then
      the inequality in~(2) of
      \cref{thm:p15-w02c-p15-001-jensen-normalization} holds for every
      admissible pair, with equality: every integrable \(f\) is zero
      a.e., and \(\varphi(f)\) is integrable only when
      \(\varphi(0)=0\). The measure with \(\mu(\{\omega_0\})=\infty\) on
      a one-point space \(\{\omega_0\}\) is an instance.
\end{enumerate}
\end{hbtheorem}

\begin{hbtheorem}{thm:p15-w02c-p15-001-jensen-extended}{proved=annex:p15-w02c-p15-001-jensen-extended}{Jensen's inequality without integrability of \(\varphi(X)\)}
\par\nobreak
Let \((\Omega,\mathcal F,\mathbb P)\) be a probability space, let
\(X\) be a real random variable with \(\mathbb E|X|<\infty\), and let
\(\varphi:\mathbb R\to\mathbb R\) be convex. Then \(\varphi(X)\) is a
random variable, \(\mathbb E[\varphi(X)^-]<\infty\), so
\(\mathbb E\varphi(X)\) exists in \((-\infty,\infty]\) in the sense of
\cref{def:p15-integrable-functions}, and
\[
 \varphi(\mathbb EX)\leq\mathbb E\varphi(X).
\]
When \(\varphi(X)\) is integrable this is the inequality of
\cref{thm:p15-jensen-integral}; otherwise
\(\mathbb E\varphi(X)=\infty\) and the inequality holds trivially.
\end{hbtheorem}

\begin{hbtheorem}{thm:p15-w02c-p15-001-jensen-integrability-limits}{proved=annex:p15-w02c-p15-001-jensen-integrability-limits}{What the two integrability hypotheses exclude}
\par\nobreak
\begin{enumerate}
\item \emph{\(X\) integrable, \(\varphi(X)\) not.} On
      \(I=\{1,2,\ldots\}\) with all subsets, let \(\mathbb P\) have the
      masses \(p(n)=3\cdot4^{-n}\) and let \(X(n)=2^n\). Then
      \(\mathbb P\) is a probability measure,
      \(\mathbb EX=3\) and \(\mathbb E(X^2)=\infty\). With
      \(\varphi(x)=x^2\), \cref{thm:p15-jensen-integral} makes no
      assertion, and \cref{thm:p15-w02c-p15-001-jensen-extended} gives
      \(9\leq\infty\).
\item \emph{\(\mathbb EX\) undefined.} On
      \(J=\mathbb Z\setminus\{0\}\) with all subsets, let \(\mathbb P\)
      have the masses \(q(m)=2^{-|m|-1}\) and let
      \(X(m)=2^m\) for \(m>0\) and \(X(m)=-2^{-m}\) for \(m<0\). Then
      \(\mathbb P\) is a probability measure and
      \(\mathbb E[X^+]=\mathbb E[X^-]=\infty\), so \(\mathbb EX\) is
      not defined and \(\varphi(\mathbb EX)\) has no meaning for any
      \(\varphi:\mathbb R\to\mathbb R\).
\end{enumerate}
\end{hbtheorem}

For a nonnegative \(X\) with \(\mathbb EX=\infty\) the expectation exists
but is not a real number, and \(\varphi\), defined on \(\mathbb R\), has
no value there; neither of Durrett's statements, nor
\cref{thm:p15-jensen-integral}, assigns one.

\begin{hbexample}{def:p15-w02c-p15-001-jensen-mass-example}{Durrett's Figure~1.6 law under a rescaled measure}
Take \(\Omega=\{1,3\}\) with all subsets, \(X(\omega)=\omega\), and
\(g(x)=x^2-3x+3\), the convex function of Durrett's Figure~1.6
\cite[p.~30]{DurrettPTEAuthor2019}; its convexity is checked in
\cref{def:p15-w02c-p15-005-jensen-direction}. The values are
\(g(1)=1\), \(g(2)=1\), \(g(3)=3\) and \(g(4)=7\). A function on
\(\Omega\) has integral \(Y(1)m_1+Y(3)m_3\) against the masses
\(m_1,m_3\), by the simple integral of
\cref{def:p15-simple-integral}.
\begin{enumerate}
\item With masses \(1/2\) and \(1/2\) (Durrett's probability measure),
      \(\int X=2\) and \(\int g(X)=2\geq1=g(2)\).
\item With masses \(1\) and \(1\) (total mass \(2\)),
      \(\int X=1+3=4\) and \(\int g(X)=1+3=4<7=g(4)\). The inequality
      reverses.
\item With masses \(1/4\) and \(1/4\) (total mass \(1/2\)),
      \(\int X=1\) and \(\int g(X)=1=g(1)\), so this \(g\) satisfies the
      inequality with equality; the constant \(\varphi\equiv1\) does not,
      since \(\int\varphi(X)=1/2<1\).
\end{enumerate}
Every function on \(\Omega\) is integrable in all three cases, so only the
total mass changes between them.
\end{hbexample}

\begin{hbwarning}{def:p15-w02c-p15-001-jensen-hypotheses-warning}{Reading the hypotheses of Jensen's inequality}
The proof of \cref{thm:p15-jensen-integral} integrates a supporting
affine function \(\varphi(c)+a(x-c)\); its integral is
\(\varphi(c)\mu(\Omega)+a(\int f\,d\mu-c\mu(\Omega))\), which reduces
to \(\varphi(c)\) at \(c=\int f\,d\mu\) when \(\mu(\Omega)=1\). This is
where the normalization enters, and
\cref{thm:p15-w02c-p15-001-jensen-normalization} shows that on a finite
measure space no other total mass can replace it. The integrability of
\(\varphi(X)\) is a convenience: by
\cref{thm:p15-w02c-p15-001-jensen-extended} the expectation of
\(\varphi(X)\) always exists, and the hypothesis only removes the case in
which it is \(\infty\). The integrability of \(X\) is what makes
\(\varphi(\mathbb EX)\) a number. For concave \(\varphi\), the
inequality reverses, by applying the theorem to the convex function
\(-\varphi\).
\end{hbwarning}

\section{Relationships}
\label{sec:p15-w02c-p15-001-relationships}

\emph{Relationships.} The proof of
\cref{thm:p15-expectation-law-integral} follows Durrett's four cases,
from indicators through simple and nonnegative functions to integrable
ones, and \cref{thm:p15-w02c-p15-001-four-step} states that scheme once.
Durrett's remark that on a space of finite measure convergence in measure
is a weaker assumption than a.e.\ convergence
\cite[sec.~1.5, p.~26]{DurrettPTEAuthor2019} concerns the two modes of
\cref{def:p15-lp-convergence}; the implication on every finite measure
space, its failure on \((\mathbb R,\lambda)\), and a sequence on the
uniform space that converges in measure and at no point are proved in
\cref{thm:p15-w02c-p15-005-ae-measure}. Case~(a) of the
notation is the one used for densities in
\cref{def:p15-probability-density} and for the conditional-density
formulas of \cref{thm:p15-conditional-density-formulas}; Fubini's
theorem in \cref{thm:p15-fubini} says when an iterated \(dx\,dy\)
integral equals the corresponding integral against \(\lambda^2\). The direct
Cauchy--Schwarz argument of \cref{thm:p15-w02c-p15-001-cauchy-schwarz}
has a conditional analogue in \cref{thm:p15-conditional-inequalities}.
The squared bound of \cref{thm:p15-w02c-p15-001-chebyshev-square} is a
special case of \cref{thm:p15-expectation-bounds}, and the interval
integrals of \cref{sec:p15-w02c-p15-001-affine-integrals} supply the
values used in \cref{thm:p15-radon-nikodym-examples}.

Durrett's other statements about the integral and expectation in his
Sections~1.5 and~1.6 are proved in
\cref{sec:p15-w02c-p15-005-integration-supplements}: monotone convergence
when the first function has an integrable negative part, his
Exercises~1.5.5 and~1.6.13 \cite[pp.~28, 37]{DurrettPTEAuthor2019}, in
\cref{thm:p15-w02c-p15-005-monotone-lower}; term-by-term integration of a
signed series with \(\sum_n\int|f_n|\,d\mu<\infty\), his
Exercise~1.5.10 \cite[p.~28]{DurrettPTEAuthor2019}, in
\cref{thm:p15-w02c-p15-005-signed-series}; the two-point reading of
Jensen's inequality and his Figure~1.6
\cite[pp.~29--30]{DurrettPTEAuthor2019} in
\cref{def:p15-w02c-p15-005-jensen-direction}, and his
Exercises~1.6.1, 1.6.2 and~1.6.12 \cite[pp.~35, 37]{DurrettPTEAuthor2019}
in \cref{thm:p15-w02c-p15-005-jensen-equality},
\cref{thm:p15-w02c-p15-005-jensen-several} and
\cref{thm:p15-w02c-p15-005-am-gm}; his display~(1.6.3)
\(\operatorname{Var}(X)\leq\mathbb EX^2\) \cite[p.~33]{DurrettPTEAuthor2019}
in \cref{thm:p15-w02c-p15-005-chebyshev-displays}; his Exercises~1.6.3,
1.6.5 and~1.6.6 \cite[pp.~35--36]{DurrettPTEAuthor2019}, which show that
the squared bound above is attained but not sharp as \(a\to\infty\), that a
mean and a variance give no positive lower bound on a tail probability,
and that \(\mathbb P(Y>0)\geq(\mathbb EY)^2/\mathbb EY^2\) for
\(Y\geq0\), in \cref{thm:p15-w02c-p15-005-chebyshev-sharp},
\cref{thm:p15-w02c-p15-005-no-lower} and
\cref{thm:p15-w02c-p15-005-positive-lower}; and his Exercise~1.6.7
\cite[p.~36]{DurrettPTEAuthor2019}, a sequence to which the
moment-control theorem \cref{thm:p15-moment-controlled-convergence}
applies although no integrable function dominates it, in
\cref{thm:p15-w02c-p15-005-no-dominator}. The same section proves his
Exercise~1.5.8 \cite[p.~28]{DurrettPTEAuthor2019}, continuity of an
indefinite integral, in \cref{thm:p15-w02c-p15-005-indefinite-continuity};
his Exercise~1.6.4 \cite[p.~36]{DurrettPTEAuthor2019}, the one-sided
Chebyshev bound, in \cref{thm:p15-w02c-p15-005-one-sided}; his
Exercises~1.6.9 and~1.6.10 \cite[pp.~36--37]{DurrettPTEAuthor2019}, the
inclusion--exclusion formula and the Bonferroni inequalities, in
\cref{thm:p15-w02c-p15-005-inclusion-exclusion} and
\cref{thm:p15-w02c-p15-005-bonferroni}; and his Exercise~1.6.14
\cite[p.~37]{DurrettPTEAuthor2019} in
\cref{thm:p15-w02c-p15-005-inverse-tail}. His Exercise~1.6.11 and
Example~1.6.12 are \cref{thm:p15-w02c-p15-001-lower-moments} and
\cref{thm:p15-w02c-p15-001-bernoulli-moments} above. The hypotheses of his
Theorems~1.5.1 and~1.6.2 \cite[pp.~25, 29]{DurrettPTEAuthor2019} are
examined in \cref{sec:p15-w02c-p15-001-jensen-hypotheses}; the
conditional form of the inequality, \cref{thm:p15-conditional-jensen},
keeps the two integrability hypotheses and replaces the single supporting
affine function by a countable family of affine minorants. Integrals on a member \(D\) of
\(\mathcal F\) against the restricted measure \(\mu|_D\) are integrals
over \(\Omega\) of zero extensions
(\cref{thm:p15-w02c-p15-003-block-integrals}); this is the identification
that the sigma-finite assembly in the proof of
\cref{thm:p15-lebesgue-decomposition} uses when it sums densities
constructed on the blocks of a partition.
